\chapter{Related Work}\label{ch:literature}

Published scores from the works below use other tasks, datasets or noise
protocols, so they serve only as context. Our result tables contain only
models we trained and evaluated ourselves.

\paragraph{CNN restoration.} DnCNN \cite{zhang2017dncnn} showed that a deep
stack of convolutions with batch normalisation learns denoising much more
easily by predicting the noise residual. FFDNet \cite{zhang2018ffdnet}
processes four half-resolution sub-images and takes a noise-level map as an
extra input, so one network covers a range of noise strengths; the author's
colour model has 852,108 parameters and reports 31.21\dB{} on CBSD68 at
$\sigma=25$. Its map describes only Gaussian noise, not blur or JPEG. Larger
transformer-based restorers reach higher quality but need tens of millions of
parameters, far outside our budget.

\paragraph{Frequency-domain restoration.} MWCNN \cite{liu2018mwcnn} replaces
pooling in a U-Net with wavelet transforms and processes subbands jointly,
which motivates our cross-band fusion. The closest work is SFNet
\cite{cui2023sfnet}, which already performs local, content-adaptive frequency
selection inside a U-shaped network. Our distinction is narrower: \fcsg{}
uses \emph{fixed} radial masks and puts the learned spatial decision in the
choice among shared $7\times7$, $5\times5$ and $3\times3$ experts, which we
can then compare with known degradation parameters.

\paragraph{Mixture of experts.} Shazeer et al.\ \cite{shazeer2017moe} route
each example to a few experts with a top-$k$ gate and balance load with
auxiliary penalties. The efficiency argument for sparse experts depends on
not running the unselected ones. Our implementation computes every expert and
mixes them with sparse weights, so sparsity here aids analysis, not speed.

\paragraph{Metrics.} PSNR is $10\log_{10}(1/\mathrm{MSE})$ on $[0,1]$ RGB
images, computed per image and averaged. SSIM compares local means, variances
and covariances in an $11\times11$ Gaussian window ($\sigma=1.5$). LPIPS
\cite{zhang2018lpips} measures distance between deep features (AlexNet
variant); lower is better, and it can disagree with PSNR.

Table~\ref{tab:gap} summarises the positioning. Each ingredient of \fcsg{} is
established; the project studies one combination of them under a tight
budget.

\begin{table}[htbp]
	\centering
	\caption{Positioning of \fcsg{} against closely related methods.}
	\label{tab:gap}
	\small
	\begin{tabular}{@{}lcccl@{}}
		\toprule
		Method & Freq.\ split & Learned selection & Spatial & Notes \\
		\midrule
		DnCNN \cite{zhang2017dncnn}      & --         & --              & --        & residual CNN \\
		FFDNet \cite{zhang2018ffdnet}    & --         & --              & --        & noise-map input \\
		MWCNN \cite{liu2018mwcnn}        & wavelet    & --              & --        & joint subbands \\
		SFNet \cite{cui2023sfnet}        & learned    & frequency       & local     & closest prior work \\
		Sparse MoE \cite{shazeer2017moe} & --         & experts         & --        & per-example routing \\
		\fcsg{} (ours)                   & radial FFT & experts, top-2  & per pixel & 97k parameters \\
		\bottomrule
	\end{tabular}
\end{table}
